\documentclass[11pt]{article}
\usepackage[utf8]{inputenc}
\usepackage[T1]{fontenc}
\usepackage{amsmath,amssymb}
\usepackage{booktabs}
\usepackage{hyperref}
\usepackage[margin=1in]{geometry}
\usepackage{enumitem}
\usepackage{tabularx}
\usepackage{xcolor}
\usepackage{parskip}

\title{\textbf{CausalCert}: Fragility-First Structural Stress-Testing for Causal DAGs}
\author{halley-labs}
\date{March 2026}

\begin{document}
\maketitle

\begin{abstract}
CausalCert is a research artifact for structural stress-testing of causal DAGs. Its executable core ranks candidate directed-pair edits by how strongly they affect d-separation, identification, and estimation, thereby surfacing the structurally sensitive parts of a causal graph. The repository also includes minimum-edit robustness-radius machinery, a large automated test suite, a quickstart pipeline, and a built-in published-DAG-style sweep. Direct execution yields nontrivial fragility rankings on both the quickstart example and the 11-example sweep, highlighting directed pairs such as \texttt{A $\to$ U2} and \texttt{Sprinkler $\to$ Rain}. These artifacts make CausalCert a practical prototype for fragility-first structural review in causal analysis and a useful foundation for further work on robustness radii.
\end{abstract}

\section{Introduction}

Applied causal analysis routinely depends on a hand-written directed acyclic graph (DAG) that encodes structural assumptions about confounding, mediation, and exclusion restrictions \cite{pearl2009,hernan2020}. Existing sensitivity tools such as the E-value and \texttt{sensemakr} ask how much \emph{unmeasured confounding} would be required to overturn a conclusion \cite{vanderweele2017,cinelli2020}; they do not say which parts of the assumed graph are most load-bearing.

That gap motivates CausalCert. The repository contains a substantial Python package and a large test suite for structural stress-testing of causal DAGs. Earlier versions of the paper tried to tell a broad story about minimum-edit robustness radii across many DAG families. The repository, however, tells a narrower and more credible story today. What actually works best is an executable scan over \emph{candidate single-edge edits}: for a treatment--outcome pair, the tool ranks directed pairs by how much changing them would perturb d-separation relations, adjustment-based identification, or effect estimation.

The central question of this recovery effort is therefore not ``what would the ideal paper claim?'' but ``what does the checked-in artifact honestly support?'' Our answer is:
\begin{enumerate}[leftmargin=*]
    \item The repository is real, runnable, and well-tested.
    \item The fragility-scoring path is executable and produces nontrivial structural diagnostics.
    \item The radius path should currently be presented as exploratory infrastructure rather than as validated evidence.
\end{enumerate}

\section{What the checked-in artifact actually computes}

The most important source-level fact is that the fragility scorer enumerates \emph{candidate single-edge edits} inside the ancestral subgraph of the treatment and outcome. For each candidate it evaluates three channels:
\begin{enumerate}[leftmargin=*]
    \item \textbf{D-separation impact:} whether the edit changes relevant conditional independences.
    \item \textbf{Identification impact:} whether the edit changes back-door identifiability.
    \item \textbf{Estimation impact:} whether the edit materially changes the estimated causal effect.
\end{enumerate}
The implementation appears in \texttt{implementation/causalcert/fragility/scorer.py} and the three channel implementations appear in \texttt{implementation/causalcert/fragility/channels.py}.

A subtle but important limitation emerges from source inspection: the public \texttt{FragilityScore} object stores only the directed pair \((u,v)\), not the \emph{edit type}. In other words, the ranking currently tells us that a particular directed pair is structurally sensitive, but not whether the high-impact perturbation was an addition, deletion, or reversal. This explains why some top-ranked outputs in the grounded runs below correspond to absent edges such as \texttt{Prenatal $\to$ Nutrition}. The present API should therefore be read as a ranking of \emph{sensitive directed pairs}, not yet as a fully typed edit explanation.

The repository also contains ILP, LP-relaxation, FPT, CDCL, and AUTO solver strategies for robustness-radius computation. Their presence is undeniable in the code. What is \emph{not} yet supported is a strong empirical claim that the end-to-end radius pipeline yields meaningful benchmark numbers under its default predicate on the checked-in synthetic runs.

\section{Implementation snapshot}

We verified the following repository facts directly:
\begin{itemize}[leftmargin=*]
    \item \textbf{167} Python source files in \texttt{implementation/causalcert}
    \item \textbf{82,007} Python source lines in that package tree
    \item \textbf{33} test files in \texttt{implementation/tests}
    \item \textbf{5} declared solver strategies (ILP, LP relaxation, FPT, CDCL, AUTO)
\end{itemize}

More importantly, we ran the full test suite rather than trusting stale paper text. The result was:
\begin{quote}
\texttt{1365 passed, 1 skipped, 13 warnings in 354.89s (0:05:54)}
\end{quote}
This does not prove that every research claim is mature, but it does show that the checked-in codebase is a real and extensively exercised artifact.

\section{Grounded experiments}

This section reports only commands that were actually executed during recovery.

\subsection{Experiment A: quickstart smoke test}

Command:
\begin{quote}
\texttt{cd implementation \&\& python3 examples/quickstart.py}
\end{quote}

The script constructs a 6-node smoking/birthweight DAG, generates synthetic data, runs the full pipeline, and prints the ranking. Table~\ref{tab:quickstart} reports the key outputs.

\begin{table}[h]
\centering
\caption{Grounded results from the checked-in quickstart script.}
\label{tab:quickstart}
\begin{tabular}{@{}ll@{}}
\toprule
Metric & Observed value \\
\midrule
Runtime & 1.57 s \\
Radius & $[0,0]$ \\
Top directed pair 1 & \texttt{Prenatal $\to$ Nutrition}, score 0.7900 \\
Top directed pair 2 & \texttt{Smoking $\to$ Nutrition}, score 0.7375 \\
Top directed pair 3 & \texttt{Nutrition $\to$ Genetics}, score 0.3150 \\
Baseline estimation & Failed with ``Unknown label type: continuous'' \\
\bottomrule
\end{tabular}
\end{table}

This run is the clearest single-piece evidence for the current story. The structural ranking is nontrivial and informative; the radius is not. The script also demonstrates an important practical limitation: the default estimation path can fail on the synthetic continuous outcomes produced by the checked-in data generator.

\subsection{Experiment B: built-in published-DAG sweep}

We repaired the existing script \texttt{implementation/examples/published\_dag\_analysis.py} so that it matches the checked-in APIs, then ran:
\begin{quote}
\texttt{cd implementation \&\& python3 examples/published\_dag\_analysis.py --samples 500 --max-nodes 8 --top-k 10}
\end{quote}

The \texttt{--max-nodes 8} restriction keeps the sweep focused on the smallest built-in DAGs and avoids claiming support for a broad multi-family benchmark. The run completed on all 11 selected DAGs. Table~\ref{tab:published} shows the exact per-DAG outputs.

\begin{table}[p]
\centering
\caption{Grounded small-DAG sweep over the built-in published-DAG library. The current pipeline surfaces nontrivial fragility rankings but returns degenerate radii and no baseline ATE on every run.}
\label{tab:published}
\small
\begin{tabularx}{\textwidth}{@{}lrrrrXrr@{}}
\toprule
DAG & $|V|$ & $|E|$ & $r_{\min}$ & $r_{\max}$ & Top directed pair & Top score & \# with score $\ge 0.4$ \\
\midrule
asia & 8 & 8 & 0 & 0 & Bronchitis $\to$ Asia & 0.6659 & 6 \\
chain & 4 & 3 & 0 & 0 & A $\to$ D & 0.3500 & 0 \\
confounded\_mediator & 6 & 8 & 0 & 0 & X $\to$ W & 0.7614 & 1 \\
diamond & 4 & 4 & 0 & 0 & M1 $\to$ M2 & 0.3500 & 0 \\
frontdoor & 4 & 4 & 0 & 0 & U $\to$ M & 0.2625 & 0 \\
hernan\_robins\_backdoor & 6 & 9 & 0 & 0 & L1 $\to$ U2 & 0.2917 & 0 \\
hernan\_robins\_mbias & 5 & 5 & 0 & 0 & A $\to$ U2 & 0.8950 & 1 \\
instrumental\_variable & 4 & 4 & 0 & 0 & Z $\to$ U & 0.3500 & 0 \\
lalonde & 8 & 13 & 0 & 0 & Age $\to$ Education & 0.0350 & 0 \\
napkin & 4 & 5 & 0 & 0 & Z $\to$ Y & 0.5250 & 1 \\
sprinkler & 5 & 5 & 0 & 0 & Sprinkler $\to$ Rain & 0.8250 & 1 \\
\bottomrule
\end{tabularx}
\vspace{0.5em}

\begin{tabular}{@{}lr@{}}
\toprule
Aggregate runtime statistic & Observed value \\
\midrule
Total time & 22.6 s \\
Mean per DAG & 2.05 s \\
Median per DAG & 0.04 s \\
Slowest DAG & lalonde (13.0 s) \\
\bottomrule
\end{tabular}
\end{table}

Three conclusions are justified by this sweep.

\paragraph{First, the fragility scan really runs.}
All 11 analyses completed and produced interpretable rankings. The highest scores were 0.8950 for \texttt{A $\to$ U2} in the M-bias example, 0.8250 for \texttt{Sprinkler $\to$ Rain}, and 0.7614 for \texttt{X $\to$ W} in the confounded-mediator example.

\paragraph{Second, the evaluation story should stay narrow.}
Only 5 of the 11 DAGs produced any directed pair with score at least 0.4, and the tool currently does not expose the edit type behind each high-scoring pair. The output is useful for structural review, but it is not yet a polished explanatory benchmark.

\paragraph{Third, the radius path is not ready to headline the paper.}
Every run in Table~\ref{tab:published} returned radius $[0,0]$, and every run reported the same baseline estimation failure on continuous outcomes. Those facts make it impossible to tell a credible paper story centered on validated minimum-edit radii.

\section{What we can truthfully claim}

The grounded evidence supports the following claims.
\begin{enumerate}[leftmargin=*]
    \item \textbf{Executable artifact claim.} CausalCert is a real, large, tested repository, not a paper-only mockup.
    \item \textbf{Fragility-first claim.} The repository already supports reproducible structural stress tests that rank sensitive directed pairs in synthetic and published-DAG-style examples.
    \item \textbf{Radius-as-infrastructure claim.} The solver stack exists in code, but the default checked-in end-to-end radius pipeline should currently be described as exploratory infrastructure rather than validated evaluation.
\end{enumerate}

Equally importantly, some claims should \emph{not} be made.
\begin{enumerate}[leftmargin=*]
    \item We should not claim a mature multi-DAG empirical benchmark for robustness radii.
    \item We should not claim that the current synthetic benchmark suite provides reliable quantitative evidence for solver quality.
    \item We should not describe the current output as a clean ranking of typed edit operations; the API currently drops the add/delete/reverse label when returning a \texttt{FragilityScore}.
\end{enumerate}

\section{Limitations and next steps}

The recovery work surfaced three concrete limitations.
\begin{enumerate}[leftmargin=*]
    \item \textbf{Default estimation mismatch.} The baseline ATE-significance predicate fails on the continuous synthetic outcomes used in the grounded runs.
    \item \textbf{Degenerate radii.} Once the baseline predicate fails, the radius collapses to $[0,0]$, which defeats the intended story.
    \item \textbf{Loss of edit type in the public ranking.} The current ranking output conflates additions, deletions, and reversals into the same directed-pair label.
\end{enumerate}

These are fixable engineering issues, but until they are fixed, the compelling paper story is the one this recovery adopts: CausalCert helps analysts \emph{stress-test structural assumptions and prioritize scrutiny}, while the minimum-edit radius machinery remains an unfinished research direction inside the same codebase.

\section{Conclusion}

After removing unsupported broad claims, the repository still supports a good paper---just a narrower one. CausalCert is best understood today as a reproducible, fragility-first structural stress-testing prototype for causal DAGs. The codebase is large and well-tested; the quickstart and small published-DAG sweep both run successfully; and the fragility rankings are nontrivial enough to be useful in practice. What the repository does \emph{not} yet support is a strong empirical story about minimum-edit robustness radii. The truthful and compelling contribution is therefore not ``we have already solved structural robustness benchmarking,'' but rather ``we have an executable tool that already helps reviewers and analysts see which directed pairs in a causal graph deserve the closest scrutiny.''

\begin{thebibliography}{9}

\bibitem{pearl2009}
J.~Pearl. \emph{Causality: Models, Reasoning, and Inference}. Cambridge University Press, 2nd edition, 2009.

\bibitem{vanderweele2017}
T.~J.~VanderWeele and P.~Ding. Sensitivity analysis in observational research: Introducing the E-value. \emph{Annals of Internal Medicine}, 167(4):268--274, 2017.

\bibitem{cinelli2020}
C.~Cinelli and C.~Hazlett. Making sense of sensitivity: Extending omitted variable bias. \emph{Journal of the Royal Statistical Society: Series B}, 82(1):39--67, 2020.

\bibitem{hernan2020}
M.~A.~Hern\'an and J.~M.~Robins. \emph{Causal Inference: What If}. Chapman \& Hall/CRC, 2020.

\end{thebibliography}

\end{document}
